\subsection*{Problema 32}

\textbf{Enunciado:} \\
Calcular la integral indefinida:
\begin{align*}
I = \int x^3 (1 + x^2)^{-1/3} \, dx
\end{align*}

\textbf{Solución:} \\
Este es un integral de un diferencial binomio de la forma $\int x^m (a + bx^n)^p \, dx$, donde:
\begin{align*}
m = 3, \quad n = 2, \quad p = -\dfrac{1}{3}
\end{align*}

Verificamos el criterio de integrabilidad calculando el valor auxiliar:
\begin{align*}
\dfrac{m + 1}{n} = \dfrac{3 + 1}{2} = \dfrac{4}{2} = 2
\end{align*}

Como el resultado es un número entero, la sustitución adecuada es igualar la expresión dentro del paréntesis a una nueva variable $u$:
\begin{align*}
u = 1 + x^2
\end{align*}

A partir de esta sustitución, despejamos el diferencial $x \, dx$:
\begin{align*}
du &= 2x \, dx \\
x \, dx &= \dfrac{1}{2} \, du
\end{align*}

También despejamos $x^2$ en términos de $u$:
\begin{align*}
x^2 = u - 1
\end{align*}

Reescribimos la integral original agrupando un factor $x$ junto al diferencial para aplicar la sustitución:
\begin{align*}
I &= \int x^2 (1 + x^2)^{-1/3} (x \, dx)
\end{align*}

Sustituimos $x^2$, $(1 + x^2)$ y $(x \, dx)$ en la integral:
\begin{align*}
I &= \int (u - 1) u^{-1/3} \left(\dfrac{1}{2} \, du\right) \\
I &= \dfrac{1}{2} \int (u - 1) u^{-1/3} \, du
\end{align*}

Multiplicamos el término $u^{-1/3}$ dentro del paréntesis aplicando leyes de exponentes:
\begin{align*}
I &= \dfrac{1}{2} \int \left( u^{2/3} - u^{-1/3} \right) du
\end{align*}

integramos cada término aplicando la regla de la potencia:
\begin{align*}
I &= \dfrac{1}{2} \left( \dfrac{u^{5/3}}{5/3} - \dfrac{u^{2/3}}{2/3} \right) + C \\
I &= \dfrac{1}{2} \left( \dfrac{3}{5} u^{5/3} - \dfrac{3}{2} u^{2/3} \right) + C
\end{align*}

Distribuimos el factor $\dfrac{1}{2}$:
\begin{align*}
I &= \dfrac{3}{10} u^{5/3} - \dfrac{3}{4} u^{2/3} + C
\end{align*}

Finalmente, regresamos a la variable original sustituyendo $u = 1 + x^2$:
\begin{align*}
I = \dfrac{3}{10} (1 + x^2)^{5/3} - \dfrac{3}{4} (1 + x^2)^{2/3} + C
\end{align*}

\begin{figure}[H]
	\centering
	\includegraphics[width=\columnwidth]{semana_05/graficas/SEM05_P32_grafica.png}
	\caption{Representación gráfica de $f(x)$ y su antiderivada $F(x)$ evaluada para la constante de integración $C = 0$ (Problema 32).}
	\label{fig:problema32}
\end{figure}
